\begin{mistake}{2026-06-11}{01}

\begin{problemcontent}
设 \(f(x)\) 在 \((-\infty, +\infty)\) 二阶可导， \(f(0) = 0\)，
\[
g(x) = \begin{cases} \frac{f(x)}{x}, & x \neq 0, \\ a, & x = 0. \end{cases}
\]
(1) 确定 \(a\) 使 \(g(x)\) 在 \((-\infty, +\infty)\) 上连续；
(2) 证明对以上确定的 \(a\)，\(g(x)\) 在 \((-\infty, +\infty)\) 上有连续一阶导数。
\end{problemcontent}

\begin{solutioncontent}
(1) 由连续性得，\(a = g(0) = \lim_{x \to 0} g(x)\)。
因为 \(f(0) = 0\)，有：
\[
\lim_{x \to 0} \frac{f(x)}{x} = \lim_{x \to 0} \frac{f(x) - f(0)}{x - 0} = f'(0)
\]
故 \(a = f'(0)\)。

(2) 当 \(x \neq 0\) 时，
\[
g'(x) = \frac{xf'(x) - f(x)}{x^2}
\]
因 \(f(x)\) 二阶可导，故 \(f'(x)\) 和 \(f(x)\) 连续，显然 \(g'(x)\) 在 \(x \neq 0\) 处连续。

当 \(x = 0\) 时，必须用导数定义求 \(g'(0)\)：
\[
g'(0) = \lim_{x \to 0} \frac{g(x) - g(0)}{x - 0} = \lim_{x \to 0} \frac{\frac{f(x)}{x} - f'(0)}{x} = \lim_{x \to 0} \frac{f(x) - xf'(0)}{x^2}
\]
使用一次洛必达法则：
\[
g'(0) = \lim_{x \to 0} \frac{f'(x) - f'(0)}{2x} = \frac{1}{2} f''(0)
\]

再求 \(g'(x)\) 在 \(x \to 0\) 的极限。因为题干仅已知 \(f(x)\) 二阶可导，未说明二阶导连续，此处使用泰勒公式最为严谨：
将 \(f(x)\) 和 \(f'(x)\) 在 \(x=0\) 处展开：
\[
f(x) = f(0) + f'(0)x + \frac{1}{2}f''(0)x^2 + o(x^2) = f'(0)x + \frac{1}{2}f''(0)x^2 + o(x^2)
\]
\[
f'(x) = f'(0) + f''(0)x + o(x)
\]
代入 \(g'(x)\) 分子：
\[
\begin{aligned}
xf'(x) - f(x) &= x[f'(0) + f''(0)x + o(x)] - [f'(0)x + \frac{1}{2}f''(0)x^2 + o(x^2)] \\
&= \frac{1}{2}f''(0)x^2 + o(x^2)
\end{aligned}
\]
因此：
\[
\lim_{x \to 0} g'(x) = \lim_{x \to 0} \frac{\frac{1}{2}f''(0)x^2 + o(x^2)}{x^2} = \frac{1}{2}f''(0)
\]
由于 \(\lim_{x \to 0} g'(x) = \frac{1}{2}f''(0) = g'(0)\)，故 \(g'(x)\) 在 \(x=0\) 处连续。
结合 \(x \neq 0\) 处的结论，\(g(x)\) 在 \((-\infty, +\infty)\) 上有连续一阶导数。
\end{solutioncontent}

\mistakeknowledge{
\begin{itemize}
  \item \textbf{易错点}：极限加减法中局部替换（如把 \(\frac{f(x)}{x}\) 提前取极限替换为 \(f'(0)\)）会丢失高阶无穷小，导致结果错误。
  \item \textbf{核心方法}：证明分段函数分界点的导数连续，必须分别求 \(g'(0)\) (用定义) 和 \(\lim_{x\to0}g'(x)\)，再判断两者是否相等。
  \item \textbf{关联知识}：仅已知“二阶可导”未说明“二阶导连续”时，求 \(\lim_{x\to0}g'(x)\) 用洛必达法则会产生逻辑漏洞，用带有皮亚诺余项的泰勒展开是最严谨的做法。
\end{itemize}}

\end{mistake}
